
\documentclass{article}
\usepackage{colortbl}
\usepackage{makecell}
\usepackage{multirow}
\usepackage{supertabular}

\begin{document}

\newcounter{utterance}

\centering \large Interaction Transcript for game `cladder', experiment `full\_v1.5\_default', episode 2927 with Qwen3.5{-}9B.
\vspace{24pt}

{ \footnotesize  \setcounter{utterance}{1}
\setlength{\tabcolsep}{0pt}
\begin{supertabular}{c@{$\;$}|p{.15\linewidth}@{}p{.15\linewidth}p{.15\linewidth}p{.15\linewidth}p{.15\linewidth}p{.15\linewidth}}
    \# & \multicolumn{2}{c}{Player} && \multicolumn{2}{c}{Game Master} \\
    \hline

    \theutterance \stepcounter{utterance}  
    & & & \multicolumn{4}{p{0.6\linewidth}}{
        \cellcolor[rgb]{0.9,0.9,0.9}{
            \makecell[{{p{\linewidth}}}]{
                \texttt{\tiny{[P1$\langle$GM]}}
                \texttt{You are an expert in causal inference. The following question is not a typical commonsense query, but rather a meticulously designed question created by a professor specializing in causal inference, intended to assess the students' mastery of the course content.} \\
\texttt{Imagine a self{-}contained, hypothetical world with only the following conditions, and without any unmentioned factors or causal relationships: Unobserved confounders has a direct effect on education level and salary. Proximity to a college has a direct effect on education level. Education level has a direct effect on salary. Unobserved confounders is unobserved. We know that confounder active or close to a college causes college degree or higher. confounder active and college degree or higher causes high salary. We observed the person lives close to a college and confounder active. Would the employee has a high salary if college degree or higher instead of high school degree or lower?} \\
\texttt{Start your answer with 'Yes' or 'No', followed by additional reasoning or evidence to support your explanation.} \\
            }
        }
    }
     \\ \\

    \theutterance \stepcounter{utterance}  
    & \multicolumn{4}{p{0.6\linewidth}}{
        \cellcolor[rgb]{0.561,0.737,0.561}{
            \makecell[{{p{\linewidth}}}]{
                \texttt{\tiny{[P1$\rangle$GM]}}
                \texttt{No} \\
\\ 
\texttt{**Reasoning:**} \\
\\ 
\texttt{In this hypothetical world, the causal structure is defined by the following relationships:} \\
\texttt{1.  **Unobserved Confounders (\$U\$)** \$\textbackslash{}rightarrow\$ **Education Level (\$E\$)**} \\
\texttt{2.  **Unobserved Confounders (\$U\$)** \$\textbackslash{}rightarrow\$ **Salary (\$S\$)**} \\
\texttt{3.  **Proximity to College (\$P\$)** \$\textbackslash{}rightarrow\$ **Education Level (\$E\$)**} \\
\texttt{4.  **Education Level (\$E\$)** \$\textbackslash{}rightarrow\$ **Salary (\$S\$)**} \\
\\ 
\texttt{The question asks for a counterfactual comparison: Would the employee have a high salary if they had a **college degree** instead of a **high school degree**, given that we observed they live close to a college (\$P=1\$) and the confounder is active (\$U=1\$)?} \\
\\ 
\texttt{To answer this, we must evaluate the counterfactual outcome \$S(E=\textbackslash{}text\{College\} | P=1, U=1)\$ versus the factual outcome \$S(E=\textbackslash{}text\{High School\} | P=1, U=1)\$.} \\
\\ 
\texttt{1.  **Analyze the Factual Scenario (\$E=\textbackslash{}text\{High School\}\$):**} \\
\texttt{    The prompt states: "confounder active and college degree or higher causes high salary." It implies that if the confounder is active but the education is *not* college degree or higher (i.e., high school), the condition for "high salary" is not met. Therefore, in the factual world where \$U=1\$ and \$E=\textbackslash{}text\{High School\}\$, the salary is **not high**.} \\
\\ 
\texttt{2.  **Analyze the Counterfactual Scenario (\$E=\textbackslash{}text\{College\}\$):**} \\
\texttt{    The prompt states: "confounder active and college degree or higher causes high salary."} \\
\texttt{    In the counterfactual scenario, we are positing that the education level is "college degree or higher" (\$E=\textbackslash{}text\{College\}\$).} \\
\texttt{    We are also given that the confounder is active (\$U=1\$).} \\
\texttt{    Since both necessary conditions for the high salary are met in this hypothetical state (\$U=1\$ AND \$E=\textbackslash{}text\{College\}\$), the causal mechanism dictates that the salary **is high**.} \\
\\ 
\texttt{3.  **Addressing the "Unobserved" Nature:**} \\
\texttt{    While \$U\$ is unobserved, the problem explicitly states we *know* the state of the confounder in this specific instance ("We observed... confounder active"). In causal inference problems of this type, if the value of a variable is known (even if it is unobserved in a general population sense), it can be used in the evaluation of the counterfactual. The "unobserved" label simply means we cannot measure it for everyone, but for this specific individual, its state is given as a premise.} \\
\\ 
\texttt{4.  **Conclusion:**} \\
\texttt{    The} \\
            }
        }
    }
    & & \\ \\

    \theutterance \stepcounter{utterance}  
    & & & \multicolumn{2}{p{0.3\linewidth}}{
        \cellcolor[rgb]{0.9,0.9,0.9}{
            \makecell[{{p{\linewidth}}}]{
                \texttt{\tiny{[GM$|$GM]}}
                \texttt{No} \\
            }
        }
    }
    & & \\ \\

    \theutterance \stepcounter{utterance}  
    & & & \multicolumn{2}{p{0.3\linewidth}}{
        \cellcolor[rgb]{0.9,0.9,0.9}{
            \makecell[{{p{\linewidth}}}]{
                \texttt{\tiny{[GM$|$GM]}}
                \texttt{yes} \\
            }
        }
    }
    & & \\ \\

    \theutterance \stepcounter{utterance}  
    & & & \multicolumn{2}{p{0.3\linewidth}}{
        \cellcolor[rgb]{0.9,0.9,0.9}{
            \makecell[{{p{\linewidth}}}]{
                \texttt{\tiny{[GM$|$GM]}}
                \texttt{game\_result = LOSE} \\
            }
        }
    }
    & & \\ \\

\end{supertabular}
}

\end{document}
